\clearpage\thispagestyle{empty}\null\clearpage
\TFMChapter{Evaluación y Experimentos Realizados}

\begin{enumerate}
\item Métricas (calidad de respuestas, latencia, coste, etc.)
\item Resultados cuantitativos y cualitativos
\item Análisis de errores y limitaciones
\end{enumerate}

\TFMSection{Introducción}
\TFMSection{Metodología de evaluación}
Definir conjunto de consultas, documentos utilizados, criterios de éxito y métricas.

\TFMSection{Evaluación del procesamiento documental}
Comprobar:
\begin{tfmbullets}
\item Calidad del Markdown generado.
\item Conservación de tablas.
\item Extracción de precios.
\item Extracción de destinos.
\item Extracción de actividades.
\item Calidad de los metadatos.
\end{tfmbullets}

\TFMSection{Evaluación de la recuperación de información}
Podéis medir:
\begin{tfmbullets}
\item Precisión de los fragmentos recuperados.
\end{tfmbullets}
\TFMTOCPageBreak

\clearpage
\begin{tfmbullets}
\item Recall.
\item Relevancia del top-k.
\item Diferencia con y sin reranking.
\item Comportamiento de filtros por destino o actividad.
\end{tfmbullets}

\TFMSection{Evaluación del sistema multiagente}
Comprobar:
\begin{tfmbullets}
\item Recopilación correcta de parámetros.
\item Selección del agente adecuado.
\item Uso correcto del RAG o API externa.
\item Transferencia de información entre agentes.
\item Finalización correcta del flujo.
\end{tfmbullets}

\TFMSection{Evaluación de las propuestas generadas}
Criterios como:
\begin{tfmbullets}
\item Coherencia del itinerario.
\item Cumplimiento del presupuesto.
\item Adaptación a fechas y preferencias.
\item Uso de información recuperada.
\item Claridad comercial.
\item Ausencia de datos inventados.
\end{tfmbullets}

\TFMSection{Comparativa de configuraciones}
Por ejemplo:
\begin{tfmbullets}
\item RAG frente a modelo sin RAG.
\item Con reranking frente a sin reranking.
\item Distintos tamaños de chunk.
\item Diferentes modelos.
\item Variaciones de prompts.
\end{tfmbullets}

\TFMSection{Resultados y análisis}

